% SPDX-License-Identifier: Apache-2.0
\section{A Program That Drives the Bus}
\label{sec:x-hostlink}

\code{Remote} sends the design's transactions out to a program that
answers them. This example runs the other way: a program on another
machine asks, and the design's bus does what it asks. The program
puts an \code{Ask} in an Ethernet frame. \code{HostLink} takes the
frame off the wire and hands the \code{Ask} on. \code{LiteMaster}
performs it as an AXI-Lite host on the design's bus, and the
\code{Answer} goes back the same way, in a frame. The two directions
share the frame of \code{RemoteLink}, at the same offsets, with two
kinds of their own, so they can share a wire and a program, and
neither takes the other's frames for its own.

This is how FireSim's host driver reaches the simulator's control
registers without PCIe. Over PCIe the driver reads and writes those
registers through the XDMA core; here the same reads and writes are
frames on the board's Ethernet port, to a machine on the other end of
a cable.

A frame can be lost, so the program sends an ask again when its
answer does not come. A write performed twice is not always harmless,
since writing some registers does something. So \code{LiteMaster}
keeps the tag and the answer of the last transaction it performed,
and an ask that repeats that tag is answered from them, without
reaching the bus.

The tags count up by one, and \code{LiteMaster} performs only the ask
whose tag follows the last one it performed. It answers any other ask
from what it kept, so the answer names the last transaction performed.
So a program need not wait for each answer before it sends the next
ask. An answer says that every ask up to its tag is done, and an
answer that names an earlier tag than the program sent says that an
ask was lost and where to start again. Each answer also names, at
byte 18 of its frame, the ask it answers, so the program can tell an
ask that came ahead of its turn, which says one was lost, from a stale
copy it sent again, which does not. FireSim's driver sends a load of
the target's memory that way, many writes before it waits. A write
whose own answer was lost could hide a failure, so the answer to a
write that was performed carries the number of transactions that have
failed so far, in place of a word. An answer from what was kept repeats
the last one, so the next answer to a write the program hears shows
the failure.

The run sends five asks, one at a time, each waiting for its answer.
It writes a word and reads it back. It asks the same read again under
the same tag, as the program does when an answer is lost, and the
count of transactions the memory saw does not move. It sends a write
whose tag skips one, as when the ask before it was lost, and the
answer names the read, the last ask performed, while the bus sees
nothing. Then it sends the skipped tag, as the program does when it
sends again from the lost ask, so the master performs it: a write
where the memory refuses, and the refusal
comes back as an error in the answer, with a count of one failure.

Both units are lowered, and the build simulates both netlists against
this run under \code{nvc} and Verilator.

\begin{figure}[htbp]
\centering
\input{hostlink_timing.tex}
\caption{An ask arriving: its bytes come in on \code{rx}, the whole
frame becomes one \code{ask}, \code{busy} rises while the read is on
the bus at \code{ar}, and the \code{back} it brings leaves as a frame
on \code{tx}.}
\label{fig:x-hostlink}
\end{figure}

\lstinputlisting[caption={\code{ex\_hostlink.rs}. Run with \code{bazel run //lib/examples:ex\_hostlink}.},label={lst:x-hostlink}]{lib/examples/ex_hostlink.rs}

\lstinputlisting[style=ascii,caption={What \code{ex\_hostlink} printed when the build ran it: each ask, its tag and answer, and how many transactions the bus had seen.},label={lst:x-hostlink-out}]{out_hostlink.txt}
